\section[Cote 21 : injectifs de $I$-torsion (proposition 2.15 et corollaire 2.16)]{Cote 21 : injectifs de $I$-torsion (proposition 2.15 et corollaire 2.16)\protect\verifie{Folder21}}\label{sec:21}

La cote 21, \emph{[Chapitre] 0 III : tapuscrit et copies de tapuscrit annotés (s.d.), notes manuscrites (s.d.)}, datée par l'inventaire \q{[vers 1960-1967]}, s'ouvre sur un tapuscrit corrigé à la main qui étudie, pour un anneau noethérien $A$ et un idéal $I$, les foncteurs contravariants sur la catégorie $\underline{C}$ des $A$-modules de type fini annulés par une puissance de $I$, et les représente par des modules $H$ dont tout élément est annulé par une puissance de $I$ (2.9 à 2.14). La proposition~2.15, énoncée page~11, caractérise ceux de ces modules qui sont injectifs ; sa démonstration, page~13, invoque le critère de Baer sous le nom de \q{Tohoku} et, pour l'étape décisive, écrit \q{En vertu de \emph{Krull}, $I^n N$ contient une intersection $N \cap I^m N$}, avec $N$ là où la ligne suivante, et l'argument, demandent $M$ (\lot{21}{1}{13}). Le corollaire~2.16 suit, puis une ligne biffée à la machine : \q{Si on rejette tout ce numéro (comme il me semble raisonnable)}.

On explicite deux notions que la lecture prend dans les numéros précédents.

\begin{definition}
Un $A$-module $H$ est \emph{de $I$-torsion} si tout élément de $H$ est annulé par une puissance de $I$. Le foncteur $\operatorname{Hom}(-, H)$ est \emph{exact à droite sur $\underline{C}$} si, pour tout $M \in \underline{C}$ et tout sous-module $N \subseteq M$, tout homomorphisme $N \to H$ se prolonge en un homomorphisme $M \to H$. Pour un $A$-module $K$, on note $H^0_I(K)$ le sous-module des éléments de $K$ annulés par une puissance de $I$.
\end{definition}

\begin{enonce}[Proposition 2.15\lie{4}]
Soient $A$ un anneau noethérien, $I$ un idéal de $A$ et $H$ un $A$-module de $I$-torsion. Pour que $H$ soit un $A$-module injectif, il faut et il suffit que le foncteur $T_H(M) = \operatorname{Hom}(M, H)$ sur la catégorie $\underline{C}$ des $A$-modules de type fini annulés par une puissance de $I$ soit exact à droite, donc exact.
\end{enonce}

\begin{enonce}[Corollaire 2.16]
Soit $K$ un $A$-module. Si $K$ est injectif, $H^0_I(K)$ l'est aussi.
\end{enonce}

\begin{lemme}\label{lem:21-puissance}
Soient $H$ un module de $I$-torsion, $M$ un module de type fini et $g : M \to H$ un homomorphisme. Il existe $n$ tel que $a\,g(x) = 0$ pour tout $a \in I^n$ et tout $x \in M$.
\end{lemme}

\begin{proof}
Soient $s_1, \dots, s_k$ des générateurs de $M$, et $e_i$ tel que $I^{e_i}$ annule $g(s_i)$. Posons $n = \max_i e_i$ ; comme $I^n \subseteq I^{e_i}$, $I^n$ annule chaque $g(s_i)$. L'ensemble des $x \in M$ tels que $I^n$ annule $g(x)$ est un sous-module (il contient $0$, est stable par somme, et par multiplication par $c \in A$ puisque $a\,g(cx) = c\,(a\,g(x))$) ; il contient les générateurs, donc c'est $M$. C'est la phrase du feuillet : l'image de $u$ est \q{un sous-module de type fini de $H$, donc annulé par une puissance de $I$}.
\end{proof}

\begin{proof}[Démonstration de la proposition 2.15]
\emph{Nécessité.} Si $H$ est injectif, tout homomorphisme d'un sous-module $N$ de $M$ vers $H$ se prolonge à $M$, pour tout $M$, a fortiori pour $M \in \underline{C}$.

\emph{Suffisance.} Supposons $\operatorname{Hom}(-, H)$ exact à droite sur $\underline{C}$. Par le critère de Baer, il suffit de montrer que, pour tout idéal $J$ de $A$, tout homomorphisme $g : J \to H$ se prolonge à $A$. L'idéal $J$ est de type fini ($A$ est noethérien) ; par le lemme~\ref{lem:21-puissance}, il existe $n$ tel que $I^n$ annule $g(J)$, et donc $g$ est nul sur $I^n J$, qui est engendré par les produits $a y$ avec $a \in I^n$, $y \in J$, et sur lesquels $g(ay) = a\,g(y) = 0$.

Par le lemme d'Artin--Rees, appliqué au module $A$, à son sous-module $J$ et à l'idéal $I$, il existe $k$ tel que $J \cap I^{p} = I^{p-k}(J \cap I^{k})$ pour tout $p \geqslant k$. Posons $m = n + k$ ; alors
\[
  J \cap I^m = I^n (J \cap I^k) \subseteq I^n J .
\]
C'est l'inclusion $N \cap I^m M \subseteq I^n N$ de la lecture, pour $M = A$ et $N = J$.

Soit $\varphi : J \to A/I^m$ la composée de l'inclusion et de la projection. Son noyau est $J \cap I^m$, contenu dans $I^n J$, sur lequel $g$ est nul. Donc $g$ se factorise par l'image de $\varphi$ : il existe $u' : \varphi(J) \to H$ avec $u'(\varphi(x)) = g(x)$ pour $x \in J$. Le module $A/I^m$ est de type fini et annulé par $I^m$, donc dans $\underline{C}$, et $\varphi(J) \simeq J/(J \cap I^m)$ en est un sous-module. Par hypothèse, $u'$ se prolonge en $w : A/I^m \to H$. La composée de la projection $A \to A/I^m$ et de $w$ prolonge $g$ : pour $x \in J$, elle envoie $x$ sur $w(\varphi(x)) = u'(\varphi(x)) = g(x)$.
\end{proof}

\begin{proof}[Démonstration du corollaire 2.16]
Le module $H^0_I(K)$ est de $I$-torsion par définition. Vérifions que $\operatorname{Hom}(-, H^0_I(K))$ est exact à droite sur $\underline{C}$. Soient $M \in \underline{C}$, annulé par $I^n$, $N \subseteq M$ et $f : N \to H^0_I(K)$. Comme $K$ est injectif, la composée de $f$ et de l'inclusion $H^0_I(K) \subseteq K$ se prolonge en $v : M \to K$. Pour $x \in M$ et $a \in I^n$, $a\,v(x) = v(ax) = v(0) = 0$ : $v$ prend ses valeurs dans $H^0_I(K)$, et fournit le prolongement cherché. La proposition~2.15 conclut.
\end{proof}

\begin{remarque}
L'inclusion d'Artin--Rees dont la démonstration a besoin est $N \cap I^m M \subseteq I^n N$, avec le module ambiant dans l'intersection, comme la lecture le rétablit ; celle de la page, $N \cap I^m N \subseteq I^n N$, est triviale pour $m \geqslant n$ et ne sert pas. Elle résulte du lemme d'Artin--Rees, non du théorème d'intersection de Krull qu'invoque la page.
\end{remarque}

\begin{remarque}
Dans~2.15, la finitude des objets de $\underline{C}$ n'intervient qu'à travers les modules monogènes $A/I^m$ : pour l'injectivité, il suffit que les homomorphismes vers $H$ se prolongent le long des inclusions $J/(J \cap I^m) \hookrightarrow A/I^m$.
\end{remarque}

\begin{remarque}
La démonstration de~2.16 n'emploie aucune finitude de $M$, seulement qu'il est annulé par une puissance de $I$ ; la noethérianité n'y sert qu'à travers~2.15.
\end{remarque}

\noindent Ne sont pas démontrés ici : les numéros~2.1 à~2.14 (foncteurs cohomologiques sur une catégorie abélienne, systèmes de générateurs, théorème d'Eilenberg--Watts, représentation d'un foncteur exact à gauche sur $\underline{C}$ par $H = \varinjlim T(A/I^n)$ en~2.9 et~2.10, critère d'annulation en bas degré en~2.12 et~2.13, premiers associés du module représentant en~2.14) ; la seconde forme de~2.15, une équivalence de catégories ; les remarques~2.17 sur les préschémas noethériens et la ind-représentabilité ; les pages~15 à~28 (conditions $(P_1)$, $(P_2)$, $(P_3)$ de prolongement à travers un fermé, pour les faisceaux d'ensembles, les catégories fibrées et les torseurs) ; la profondeur le long d'un fermé et les énoncés de prolongement de type Hartogs.
